% Formal annotations.
%
% Each theorem-like environment carries a machine-readable provenance record.
% The six macros are typographically empty and take comma-separated identifiers:
%
%   \coverage{...}  absent, fragment, conditional_deduction, or complete;
%   \lean{...}      reviewer terminals, without the common namespace;
%   \uses{...}      manuscript dependencies: conceptual in a statement,
%                   logical at the end of its proof;
%   \proves{...}    the label established by a detached proof;
%   \imports{...}   external-result identifiers resolved by the source registry;
%   \evidence{...}  computational bundles resolved by the evidence registry.
%
% Details live in lean/verification/claims.json, verification/imported-sources.json,
% and verification/evidence.json. The source-only checker resolves every identifier
% and checks the declared dependency graph and reviewed statement/source-signature
% digests. Hashing does not elaborate Lean. An absent statement carries no \lean.
% Proof annotations are grouped at the end, before \end{proof}, so the empty
% macros do not change the run-in proof heading. The present paper has no theorem
% whose proof rests on computational evidence; its evidence registry says so.
% The uses/lean/proves vocabulary follows the Lean blueprint conventions; the
% conceptual/logical distinction follows KnowTeX. All macros take identifiers only.
\newcommand{\lean}[1]{}
\newcommand{\coverage}[1]{}
\newcommand{\uses}[1]{}
\newcommand{\imports}[1]{}
\newcommand{\evidence}[1]{}
\newcommand{\proves}[1]{}
